% generated by src/make_tables.py -- do not edit
\begin{table}[!t]
\centering
\caption{Celah optimalitas menurut \eqref{eq:gap} pada instans yang $\chi(G)$-nya terbukti. Setiap sel berisi $G_{\mathrm{WP}}/G_{\mathrm{DS}}$ dalam persen; tanda -- berarti sel itu di luar cakupan eksak.}
\label{tab:optimality}
\footnotesize
\setlength{\tabcolsep}{3pt}
\begin{tabular}{@{}lcccc@{}}
\toprule
 & $p=0.1$ & $p=0.3$ & $p=0.5$ & $p=0.7$ \\
\midrule
\multicolumn{5}{@{}l}{\emph{Eksperimen A (acak)}} \\
$|V|=20$ & 10.56/1.11 & 9.17/4.17 & 11.37/3.56 & 7.53/1.73 \\
$|V|=40$ & 23.33/7.50 & 20.78/11.67 & 25.19/14.58 & 17.75/12.21 \\
$|V|=60$ & 22.50/3.33 & 30.89/17.14 & 28.09/15.91 & -- \\
$|V|=80$ & 26.33/11.67 & -- & -- & -- \\
\midrule
\multicolumn{5}{@{}l}{\emph{Eksperimen B (pendaftaran)}} \\
kecil & \multicolumn{4}{c}{0.00/0.00} \\
menengah & \multicolumn{4}{c}{0.78/0.25} \\
besar & \multicolumn{4}{c}{0.45/0.47} \\
\midrule
Total & \multicolumn{4}{c}{15.75/7.06} \\
\bottomrule
\end{tabular}
\end{table}
